\subsection{Controlled transitions and regret-scaled relaxation}\label{sec:controlled-audit}
The controlled audit uses a condition regime $i=0,\ldots,n-1$, three or four actions, and reward $i/(n-1)-c_a$ with a small rational price perturbation. More intensive actions shift transition probability from deterioration toward improved condition and replacement. The installed rule, nonnegative revision weights, uniform initial regime law, and terminal reward are fixed by a complete rational primitive contract. Thus the transition law is genuinely endogenous, unlike the exogenous service economies below. The primitives are designed, not estimated.

Eight cases were declared before evaluation, together with solver and verifier source hashes. During the initial execution, accumulating rational expressions with unrelated denominators made certificate arithmetic unnecessarily expensive; an integer-string limit also interrupted a run. The replacement accumulates outward dyadic enclosures of individual summands. Original source, results, the amendment, and both remote hash registrations remain in the computation record. No case, operating allowance, action set, transition generator, or LP proposal rule was changed. This is a predeclared-case experiment with a disclosed arithmetic amendment, not an untouched code-frozen holdout.

Both methods use the reference-eliminated Bellman formulation, identical operating-regret bounds and probability caps, the same LP engine and tolerances, and the same exact candidate repair. Method U uses ordinary implementation-cost bounds. Method S additionally uses the savings cones and contracted savings bounds. This is a same-object root-relaxation ablation; neither method branches in this experiment. The 30-second cap applies to the LP call, not to model construction or proof checking. All those costs are recorded separately. A stopped LP contributes a sound fallback bound and the operating-optimal feasible reference, not an unverified primal or dual estimate.

\begin{table}[t]
\caption{All predeclared controlled cases: root-relaxation certificates}\label{tab:controlled}
\centering\small
\begin{tabular}{rrrrlrrrl}
\toprule Case&$n$&$T$&$m$&Method&Lower&Upper&Gap/upper (\%)&LP status\\\midrule
\tablerows{revisions/2026-09-25-r43/paper/generated/controlled_rows.tex}
\bottomrule
\end{tabular}
\par\smallskip\footnotesize U: unscaled. S: regret-scaled. Endpoints are outward decimal displays of independently checked rational endpoints. ``Cap'' retains the complete sound fallback interval. These are root bounds for the same constrained common-policy optimum, not floating-point LP objective intervals. The supplement reports all allowances, dimensions, timings, memory, proof sizes, and worst-case constants.
\end{table}

The savings bounds materially reduce the certified interval in \RControlledImproved{} of the eight pairs (Table~\ref{tab:controlled}); a reduction below $10^{-8}$ in absolute width is not counted. The eight-regime, 32-period, four-action case at allowance $10^{-3}$ gives $[\RLongLower,\RLongUpper]$, a relative width of \RLongRelative\%. Other cases receive little benefit, and \RControlledCaps{} LP calls reach their cap. In particular, a stronger mathematical relaxation need not give a stronger \emph{computed} certificate within a fixed time limit. Intersecting any independently valid intervals is permitted, but the table retains the individual methods rather than hiding an unfavorable run behind a portfolio.

\begin{table}[t]
\caption{Three orders of allowance refinement in one fixed eight-period model}\label{tab:allowance}
\centering\small
\begin{tabular}{rrrrrr}
\toprule $\varepsilon$&U width&S width&S width$/\varepsilon^2$&S total (s)&S check (s)\\\midrule
\tablerows{revisions/2026-09-25-r43/paper/generated/scaling_rows.tex}
\bottomrule
\end{tabular}
\par\smallskip\footnotesize The model has three regimes and three actions. U and S are defined in Table~\ref{tab:controlled}. Total time includes primitive construction, LP proposal, and constructor arithmetic; the independent check is additional. Widths are displayed upward. The ratio is descriptive: the theorem is a fixed-model upper rate, not an assertion that every pair of observed widths has slope two.
\end{table}

At allowance $10^{-4}$ the scaled positive-cost interval has width at most \RScaleGap{}, compared with \RScaleUnscaled{} for the unscaled relaxation. This calculation illustrates why tight allowances can help a separated-action formulation even though they worsen the generic repair denominator. It does not imply that the action-separation constants are uniformly favorable. The supplement reports $\underline d$, the savings-cone slopes, and each term of \eqref{eq:quadratic-width}; these sufficient bounds can be much larger than the measured interval. None of the eight larger-state cases is relabeled as having met an absolute target merely because its relative gap is small.

\subsection{An executed controlled atomic continuum certificate}\label{sec:fiber-execution}
Consider a continuum of observable establishment sizes $z\in[0,1]$, fixed over the decision horizon. Condition evolves according to the four-regime, three-action controlled maintenance law specified in the replication contract; the horizon is eight and $\beta=0.95$. Operating rewards and the terminal reward are multiplied by $s(z)=1+z$, while revision costs are multiplied by $w(z)=1+z/2$. Size and the initially uniform regime are independent, and size is uniform. The operating allowance is $0.01$ for every size, condition, and restart. Every transition is atomic in the continuous size coordinate. Actions nonetheless change the condition transition probabilities.

Writing $v_0$ for the base finite-model optimum, Theorem~\ref{thm:fibers} gives
\begin{equation}\label{eq:scaled-fiber}
 \KR=\int_0^1(1+z/2)\,v_0\!\left(\frac{0.01}{1+z}\right)dz.
\end{equation}
This is an optimization identity, not an assumed restriction to a parametric policy. On a cell $[a,b]$, monotonicity of $v_0$ in its allowance supplies a lower endpoint from the finite solve at $0.01/(1+a)$ and a feasible upper policy from the tighter solve at $0.01/(1+b)$. Multiply these endpoints by the exact positive weight $(b-a)+(b^2-a^2)/4$ and sum. The upper policy uses the observed size cell at every date and retains common continuations.

\begin{table}[t]
\caption{Controlled atomic continuum: paired partition bounds}\label{tab:controlled-continuum}
\centering\small
\begin{tabular}{rrrrrr}
\toprule Cells&Lower&Upper&Gap/upper (\%)&Finite widths&Partition remainder\\\midrule
\tablerows{revisions/2026-09-25-r43/paper/generated/fiber_rows.tex}
\bottomrule
\end{tabular}
\par\smallskip\footnotesize All 17 finite endpoint problems and all four weighted sums are independently checked. ``Finite widths'' is the sum of the upper-endpoint fiber widths with exact cell weights. The remainder is the displayed total interval width minus that quantity; it includes value change and differences between adjacent lower endpoints, not an independently estimated quadrature error.
\end{table}

The 16-cell interval is $[\RFiberLower,\RFiberUpper]$, with relative width \RFiberRelative\%. It is a paired global certificate for an uncountable-state controlled atomic model, not a density approximation or a collection of pointwise policy tests. Partition refinement reduces the displayed interval, but the finite root widths remain material. Arbitrarily refining size cells alone does not force this execution to reach zero width: the finite fibers must also be solved more accurately, as required by Proposition~\ref{prop:fiber-cells}. The original action-dependent continuous-condition maintenance models are different models and retain their existing status.
